\section{Discussion}
\label{sec:discussion}

\paragraph{Ranked hardening guidance.}
Our empirical results demonstrate a clear hierarchy of defensive efficacy. Userspace kernel isolation (\textbf{B7}, gVisor) provides the single strongest defense (ARR $=0.97$), eliminating escapes across device passthrough (\texttt{privileged}), socket mounting (\texttt{docker\_socket}), and process namespace sharing (\texttt{pid\_namespace}). User-namespace remapping (\textbf{B5}) is a partial second-line control. In the reference-exploit check it closes host-filesystem write paths (\texttt{hostpath\_etc}), where gVisor is ineffective, and the capability-limited scenarios. It leaves \texttt{docker\_socket} OPEN (the reference exploit succeeds under B5), and it does not apply to \texttt{privileged} or \texttt{pid\_namespace} containers, because the Docker daemon rejects those containers when user namespaces are enabled. In contrast, standard container-level flags (\textbf{B1--B4}) provide minimal to zero protection against structural misconfigurations. We recommend a prioritized hardening hierarchy: (1) enforce userspace kernels (gVisor) for container workloads, (2) enable daemon-level user-namespace remapping (B5) where the workload permits it, recognizing that it does not cover the \texttt{docker\_socket} path, and (3) apply strict seccomp profiles (B3) to restrict syscall surfaces.

\paragraph{Real security versus model friction.}
A central finding of our work is that empirical LLM escape rates alone can produce misleading security assurances. On capability-limited scenarios (\texttt{sys\_admin}, \texttt{sys\_module}, \texttt{dac\_read\_search}), the LLM scored $0\%$ escape success across all controls---including the unhardened baseline (B0). Without mechanistic verification, an evaluator might falsely conclude that capability dropping (B1) or read-only filesystems (B4) successfully defended the container. RQ2 reference exploit cross-checks revealed that B1 and B2 left the underlying capabilities fully open, and B4 merely introduced script compilation friction in \texttt{/tmp}. Security controls must be validated against LLM-independent reference exploits to confirm true architectural boundaries rather than temporary attacker friction.

\paragraph{Defense-in-depth.}
The composite control (\textbf{B8}), which bundles capability dropping (B1), no-new-privileges (B2), strict seccomp (B3), and read-only rootfs (B4), proved largely redundant relative to its constituent parts. On misconfiguration scenarios (\texttt{privileged}, \texttt{docker\_socket}, \texttt{pid\_namespace}), B8 achieved SR $=0.90$--$1.00$, demonstrating that bundling weak container-level controls fails to produce cumulative security when the primary attack vector (e.g., host block device or host daemon socket) remains exposed. Defense-in-depth is effective only when controls address orthogonal attack vectors rather than layering ineffective constraints on the same unhardened boundary.

\paragraph{Relation to the ``ceiling effect'' hypothesis.}
Our findings refine the ``ceiling effect'' hypothesis proposed in web exploitation benchmarks such as PACEbench~\cite{pacebench}. In web applications, production WAFs held against all seven frontier models tested in PACEbench. In container environments, however, standard container-level hardening (B1--B4) fails to induce a ceiling effect against structural misconfigurations: LLM agents successfully navigate hardened environments to retrieve host flags. A true defensive ceiling in container sandboxes requires fundamental isolation primitives---namely userspace kernel virtualization (gVisor) or namespace remapping (B5).
